\documentclass[11pt]{article}
\usepackage[margin=1in]{geometry}
\usepackage{amsmath, amssymb, amsthm}
\usepackage{hyperref}

\newtheorem{theorem}{Theorem}
\newtheorem{definition}{Definition}
\newtheorem{proposition}{Proposition}

\title{Combinatorial Auctions: Winner Determination, VCG, and Core Pricing}
\author{Abhi Wadhwa}
\date{}

\begin{document}
\maketitle

\begin{abstract}
Notes on the theory behind combinatorial auctions. I cover the winner determination problem (WDP), the VCG mechanism, and core pricing. The main thread is the tension between efficiency, truthfulness, and revenue---VCG gives you the first two but can fail badly on the third, and core pricing is the attempt to fix that. Along the way, I try to build intuition for why these problems are hard and why the standard solutions work (and don't).
\end{abstract}

\section{Why Combinatorial?}

In a standard auction, you sell items one at a time. But what if the items interact? Maybe a TV license for New York is worth \$10M on its own, but the license for New York \emph{plus} New Jersey together is worth \$30M because you can offer seamless coverage. Or maybe two paintings by the same artist are substitutes---having both isn't much better than having one.

These are \textbf{complementarities} and \textbf{substitutabilities}, and they're everywhere in practice: spectrum licenses, airport landing slots, procurement contracts, ad auctions. The problem with selling items separately is that bidders can't express these synergies. They might overbid on individual items hoping to win the bundle, or underbid because they're afraid of getting stuck with half a complement.

A \textbf{combinatorial auction} lets bidders bid on \emph{bundles} (subsets) of items directly. This is more expressive, but it creates a much harder optimization problem.

\section{The Winner Determination Problem}

We have a set of items $M = \{1, \ldots, m\}$ and $n$ bids, where bid $j$ is for bundle $S_j \subseteq M$ at price $v_j \geq 0$. The \textbf{winner determination problem} (WDP) selects a subset of bids to maximize total value, subject to the constraint that each item is allocated at most once:
\begin{align}
\max \quad & \sum_{j=1}^{n} v_j \cdot x_j \label{eq:wdp-obj} \\
\text{s.t.} \quad & \sum_{j : i \in S_j} x_j \leq 1 \quad \forall \, i \in M, \label{eq:wdp-item} \\
& x_j \in \{0, 1\} \quad \forall \, j. \label{eq:wdp-binary}
\end{align}

This is a \textbf{weighted set packing} problem, and it's NP-hard. The reduction is straightforward: weighted set packing is a well-known NP-hard problem, and WDP is literally an instance of it.

That said, ``NP-hard'' doesn't mean ``impossible in practice.'' Modern ILP solvers (like CBC, CPLEX, or Gurobi) handle real-world WDP instances with thousands of bids quite efficiently. The LP relaxation (replacing $x_j \in \{0,1\}$ with $0 \leq x_j \leq 1$) often has integral solutions, especially when there isn't too much competition for items. And even when it doesn't, branch-and-bound closes the gap quickly for most practical instances.

Notice that the WDP objective \eqref{eq:wdp-obj} maximizes \textbf{social welfare}---the total value to all winning bidders. This is the ``efficient'' allocation, the one that creates the most surplus. But we're assuming the bid values $v_j$ are truthful. If bidders can lie about their values, we need a mechanism that incentivizes honesty.

\section{The VCG Mechanism}

The \textbf{Vickrey-Clarke-Groves} mechanism is the standard tool for getting truthful bids in auctions. The idea goes back to Vickrey's second-price auction: make each bidder pay not their bid, but the \textbf{externality} they impose on others.

Here's how it works in the combinatorial setting. Solve the WDP to get the efficient allocation $(S_1^*, S_2^*, \ldots)$ with total welfare $V^*(N)$. For each winning bidder $i$, solve the WDP \emph{again} with bidder $i$ excluded, getting welfare $V^*(N \setminus \{i\})$. Bidder $i$'s VCG payment is:
\[
p_i = V^*(N \setminus \{i\}) - \left[V^*(N) - v_i(S_i^*)\right].
\]

In words: $p_i$ equals the total value that others would have gotten without $i$, minus the total value others actually get with $i$. This is the ``damage'' that $i$'s presence causes to the other bidders.

The VCG mechanism has three beautiful properties:

\textbf{Truthfulness (DSIC).} Reporting your true valuation is a dominant strategy. Here's why: your payment doesn't depend on your own bid (only on others' values and the structure of the efficient allocation). Your bid only affects \emph{which bundle you get}. Since the mechanism allocates efficiently, bidding truthfully ensures you get the bundle that maximizes your surplus $v_i(S_i^*) - p_i$. Any deviation either has no effect or gives you a worse bundle.

\textbf{Individual rationality.} Every winner pays at most their value: $p_i \leq v_i(S_i^*)$. This follows because $V^*(N) \geq V^*(N \setminus \{i\})$---adding a bidder can't decrease the optimal welfare.

\textbf{Efficiency.} The allocation maximizes total value, by construction.

Computing VCG payments requires solving $n + 1$ ILPs: one for the full problem and one per winning bidder with that bidder excluded. This is the computational cost of truthfulness.

\section{The Problem with VCG: Revenue and the Core}

So VCG is truthful, individually rational, and efficient. What could go wrong?

\textbf{Revenue.} VCG payments can be disturbingly low. Consider two bidders who each want the same single item, bidding \$100 and \$99. VCG gives the item to the first bidder at a price of \$99 (second-price logic). Fine. But now consider two bidders with complementary valuations: Alice values $\{A, B\}$ at \$100 but each item alone at \$0, and Bob values $\{A, B\}$ at \$99 but each item alone at \$0. VCG gives the bundle to Alice at a price of \$99. Still fine.

But things get weird with more complex complementarities. There are examples where VCG revenue is \emph{zero} even though the items are highly valuable---essentially because no single bidder's exclusion changes the outcome.

\textbf{Core violations.} The \textbf{core} of the auction is the set of payment vectors where no coalition of bidders and the auctioneer can profitably deviate. A payment vector $p$ is in the core if for every coalition $C \subseteq N$:
\[
\sum_{i \in C} p_i \geq V^*(C) - \left[V^*(N) - \sum_{i \in C} v_i(S_i^*)\right].
\]

This says the payments from the bidders outside $C$ must be high enough that the auctioneer wouldn't prefer to sell to $C$ instead. If VCG payments violate this, the auctioneer has a legitimate complaint: ``I could have gotten more by dealing with a different set of bidders.''

And indeed, \textbf{VCG payments are not always in the core}. The classic example is the ``threshold problem'': three items, three bidders, each wanting a pair. The complementarities create a situation where VCG payments are too low because each individual bidder's exclusion doesn't change things (the other two cover all items), but the auctioneer collectively gets very little.

\section{Core-Selecting Auctions}

Core-selecting auctions fix the VCG revenue problem by finding payments in the core. The idea is: keep the efficient VCG allocation, but adjust the payments upward to lie in the core. Among all core-compatible payment vectors, pick one that minimizes some objective (e.g., minimize total payment, or minimize the distance from VCG payments).

The core constraints form a polyhedron, and finding the minimum-revenue core point is an LP. The FCC's combinatorial clock auction, used for spectrum license sales, is essentially a core-selecting auction.

The tradeoff is that \textbf{core-selecting auctions sacrifice truthfulness}. Once you're no longer using VCG payments, bidders may have an incentive to shade their bids. The argument for core-selecting auctions is pragmatic: the revenue loss from VCG is unacceptable in practice, and the degree of strategic manipulation in core-selecting auctions is empirically small.

I find the whole VCG-vs-core debate fascinating because it illustrates a fundamental impossibility: \textbf{you can't have efficiency, truthfulness, individual rationality, and core stability simultaneously} (in general). Something has to give. VCG gives up core stability; core-selecting auctions give up truthfulness. Pick your poison.

\section{Computational Remarks}

The WDP is the computational bottleneck. Everything else (VCG payments, core checking) reduces to solving additional WDP instances. So the practical question is: how hard is the WDP for realistic inputs?

For structured inputs (single-minded bidders, items on a line, hierarchical bundles), there are polynomial-time algorithms or strong LP relaxations. For general inputs, you're relying on ILP solvers. The good news is that the LP relaxation of WDP is often very tight, and modern branch-and-bound can handle instances with thousands of bids.

Checking whether VCG payments are in the core requires verifying the core constraint for every coalition $C$, which is $2^n$ constraints. In practice, you can use separation: start with VCG payments, check if any coalition violates the core constraint (this is itself an optimization problem), and if so, you know VCG isn't in the core. Finding the minimum-revenue core point can also be formulated as an LP with separation.

\begin{thebibliography}{9}
\bibitem{vickrey1961} W.~Vickrey, ``Counterspeculation, auctions, and competitive sealed tenders,'' \textit{Journal of Finance}, vol.~16, no.~1, pp.~8--37, 1961.

\bibitem{clarke1971} E.~H. Clarke, ``Multipart pricing of public goods,'' \textit{Public Choice}, vol.~11, pp.~17--33, 1971.

\bibitem{groves1973} T.~Groves, ``Incentives in teams,'' \textit{Econometrica}, vol.~41, no.~4, pp.~617--631, 1973.

\bibitem{cramton2006} P.~Cramton, Y.~Shoham, and R.~Steinberg, \textit{Combinatorial Auctions}, MIT Press, 2006.

\bibitem{devries2003} S.~de~Vries and R.~V. Vohra, ``Combinatorial auctions: A survey,'' \textit{INFORMS Journal on Computing}, vol.~15, no.~3, pp.~284--309, 2003.

\bibitem{day2008} R.~Day and P.~Milgrom, ``Core-selecting package auctions,'' \textit{International Journal of Game Theory}, vol.~36, no.~3--4, pp.~393--407, 2008.
\end{thebibliography}

\end{document}
